\documentclass[12pt]{article}
\usepackage[utf8]{inputenc}
\usepackage[T1]{fontenc}
\usepackage{hyperref}
\title{Tzotzil, Chamula Language Resources for K-12 Curriculum Development}
\author{Compiled by Claude Sonnet 5 (claude-sonnet-5) for Andr\'as Kornai}
\date{}
\begin{document}
\maketitle

\section{Introduction}
Chamula Tzotzil (ISO 639-3 tzc) is a Mayan language of Chiapas, Mexico. The identification given with the task lists the family as ``Bookkeeping'', which looks like a data error. The searches found no dictionary or grammar specific to Chamula, and most substantial Tzotzil lexical work covers Zinacant\'an.

\section{Primary Educational Resources}
A 1970 ILV Spanish and Chamula Tzotzil phrase booklet of 4 pages is recorded in the Colegio de M\'exico repository. I did not confirm its URL, so it is not listed with one. According to the Wikipedia article, some primary schools and a secondary school in central Chiapas teach in Tzotzil. Whether they use Chamula materials is not stated.

\section{Linguistic and Reference Materials}
\begin{itemize}
\item Tzotzil language, Wikipedia: \url{https://en.wikipedia.org/wiki/Tzotzil_language}
\item Laughlin's 1975 Smithsonian dictionary of Zinacant\'an Tzotzil (a different variety, useful only for comparison): \url{https://repository.si.edu/handle/10088/1367?show=full}
\item Endangered Languages Project entry, which lists \emph{Norma de escritura de la lengua tsotsil} (2011); I could not see whether it covers Chamula: \url{https://endangeredlanguages.com/node/106949}
\item William O. Bright Papers, Tzotzil materials including an SIL Tzotzil-Spanish-Tzotzil dictionary of 1949 (dialect not stated): \url{https://indigenousguide.amphilsoc.org/entry/39182}
\end{itemize}

\section{Religious and Community Resources}
None were found in the searches.

\section{Digital and Audio Resources}
\begin{itemize}
\item Living Dictionary for Tzotzil (Bats'i K'op); the dialect was not confirmed: \url{https://livingdictionaries.app/batsi-kop-tsotsil-tsot/entry/uvOCnOqcSqlFkntYgS1u}
\end{itemize}

\section{Curriculum Development Context}
Little Chamula-specific material turned up. A curriculum would likely have to adapt Tzotzil resources from other varieties and develop new Chamula materials with the community.
\end{document}